\naslov{Osnovni tipi PDE}

Uporabljamo notacijo s predavanj:
\begin{itemize}
\item $\alpha = (\alpha_1, \ldots, \alpha_n) \in \N_0^n$ je multiindeks,
\item $\abs{\alpha} = \alpha_1 + \cdots + \alpha_n$,
\item odvod funkcije $u$ z multiindeksom $\alpha$ je
  \[
	D^\alpha u = \frac{\partial^{\abs{\alpha}}}{\partial x_1^{\alpha_1} \ldots
	  \partial x_n^{\alpha_n}},
  \]
\item za $k \in \N$ označimo skupek odvodov $k$-tega reda kot $D^k u =
  \{D^\alpha u \such \abs{\alpha} = k\}$.
  To včasih obravnavamo kot množico, včasih pa kot vektor ali matriko.
\end{itemize}

\begin{definicija}
  \pojem{Parcialna diferencialna enačba} je enačba, ki vsebuje neznano funkcijo
  $u$ vsaj dveh spremenljivk ter nekatere njene parcialne odvode.
\end{definicija}

\begin{definicija}
  \pojem{Parcialna enačba $k$-tega reda} je enačba oblike
  \[
	F(x, u(x), D u(x), \ldots, D^k u(x)) = 0,
  \]
  kjer je $x \in U\odpp\R^n$ in $F: U \times R \times \R^n \times \cdots \times
  R^{n^k} \to \R$ dana.
\end{definicija}

Rešitev PDE je vsaka funkcija, ki enačbi zadošča.
Lahko je klasična ali posplošena rešitev; za klasično obstajajo vsi odvodi $Du,
\ldots, D^k u$, ki jih vstavimo v $F$, posplošena rešitev pa je rešitev v smislu
integracije po delih.

\begin{primer}
  Dana je enačba $\Delta u = f$.
  Če je $u \in \zvodvi{2}$ in označimo $\Delta u = g$, za vsak $\phi \in
  \zvodv{\infty}{U}$ s kompaktnim nosilcem velja
  \[
	\sk{\Delta u, \phi} = \sk{g, \phi} = \int_U g(x) \phi(x) dx.
  \]
  Iz integracije po delih sledi $\sk{\Delta u, \phi} = \sk{u, \Delta \phi}$
  oziroma $\sk{u, \Delta \phi} = \sk{g, \phi}$.
  Tu zahtevamo samo, da je $u$ dovolj regularen, da lahko izračunamo integral v
  skalarnem produktu.
  Šibka rešitev je vsaka $u \in \zvezne{U}$, ki zadošča $\sk{u, \Delta \phi} =
  \sk{f, \phi}$ za vsak $\phi$.
\end{primer}

Poznamo štiri osnovne tipe parcialnih diferencialnih enačb.
\begin{itemize}
\item Linearna PDE je linearna kombinacija parcialnih odvodov
  \[
	\sum_{\abs{\alpha} \le k} a_\alpha(x) D^\alpha u(x) = f(x).
  \]
\item Semilinearna PDE je linearna samo v odvodu najvišjega reda
  \[
	\sum_{\abs{\alpha} = k} a_\alpha(x) D^\alpha u + b(x, u, Du, D^{k-1} u) = 0.
  \]
\item Kvazilinearna PDE je enačba oblike
  \[
	\sum_{\abs{\alpha} = k} a_\alpha(x, u, Du, \ldots, D^{k-1}u) \cdot D^\alpha
	u + b(x, u, \ldots, D^{k-1} u) = 0.
  \]
\item Povsem nelinearna PDE je enačba, ki je nelinearno odvisna od odvodov
  najvišjega reda.
\end{itemize}

\vprasanje{Definiraj linearne, semilinearne in kvazilinearne PDE.}

\begin{definicija}
  Pravimo, da je problem PDE z začetnim ali robnim pogojem \pojem{dobro
	postavljen}, če ima naslednje lastnosti:
  \begin{itemize}
  \item rešitev obstaja,
  \item rešitev je enolična,
  \item rešitev je zvezno odvisna od začetnih pogojev.
  \end{itemize}
\end{definicija}

% LocalWords:  multiindeks multiindeksom
